\documentclass[10pt,aspectratio=169,utf8]{beamer}
% \documentclass[14pt,notes only, handout,aspectratio=169,utf8]{beamer}
% \documentclass[10pt,handout,aspectratio=169,utf8]{beamer}

  % We need this package to properly import our dependencies below.
  \usepackage{import}  % subimport
  \subimport{style/}{all.tex}

  \title{Clean Test}
  \subtitle{A modern testing framework}
  \author{Philipp Ochsendorf}
  \date{November 17th, 2022}

\begin{document}

\begin{frame}
 \maketitle
\end{frame}

\begin{frame}
  \frametitle{DeepL: We are hiring}

  \begin{columns}[onlytextwidth]
    \begin{column}{.4\textwidth}
      \begin{center}
        \href{%
          https://www.deepl.com/translator
        }{%
          \includegraphics[width=\textwidth]{img/deepl}
        }

        \vspace{15pt}

        \href{%
          https://jobs.deepl.com/
        }{
          \Large{We are hiring!}
        }
      \end{center}
    \end{column}

    \begin{column}{.55\textwidth}
      \begin{itemize}
        \item
        Meaningful work: \\
        we break down language barriers

        \item
        Fast growing European startup

        \item
        Interesting challenges at scale

        \item
        Modern \texttt{C++}, latest tools

        \item
        Work independently with smart colleagues

        \vspace{10pt}

        \item Many open positions in engineering teams, \\
        several with a C++ focus
        (\href{https://jobs.deepl.com/o/senior-backendentwickler-c-wmd-deremote-koln-oder-paderborn}{backend},
        \href{https://jobs.deepl.com/o/senior-data-engineer-wmd-language-data-3}{research})
      \end{itemize}
    \end{column}
  \end{columns}
\end{frame}
\note{
  We develop neural networks to help people work with language.
}

\section{Introduction}
\begin{frame}[fragile]
  \frametitle{Clean Test: a modern testing framework}

  \begin{columns}[T,onlytextwidth]
    \begin{column}{.55\textwidth}
      \begin{itemize}
        \item Standard \texttt{C++}20, no macros
        \item Parallel test execution, \\ thread safe expectations
        \item CI ready: JUnit export, UTF-8 support
        \item \texttt{CMake} library without dependencies
        \item Liberal \texttt{BSL-1.0} license
      \end{itemize}
    \end{column}
    \begin{column}{.45\textwidth}
      \includeDemo{Overview.cpp}
    \end{column}
  \end{columns}
\end{frame}
\note{
  \begin{itemize}
    \item
    \textbf{Nowadays} testing sequentially wastes resources; simulataneous failures are essential.

    \item
    Not header only: easy to integrate with \texttt{CMake}, no need to recompile the framework over and over again.

    \item
    \code{ct::main} not necessary

    \item
    no macros $\Rightarrow$ experienced developers can use IDE to understand example

    \item
    3 Concepts: suites, tests and expectations.
  \end{itemize}
}

\section{Suites and Tests}
\begin{frame}
 \frametitle{Suites and Tests}

 \begin{columns}[onlytextwidth]
  \begin{column}{.45\textwidth}
    \begin{overlayarea}{\textwidth}{.9\textheight}
      \only<1 | handout:0>{
        \includeDemo[%
          linerange = {1-12},%
        ]{Intro.cpp}
      }
      \only<2 | handout:0>{
        \includeDemo[%
          linerange = {1-17,22},%
        ]{Intro.cpp}
      }
      \only<3->{
        \includeDemo{Intro.cpp}
      }
    \end{overlayarea}
  \end{column}

  \begin{column}{.55\textwidth}
    \begin{description}
      \item[Suite] test initialization executed automatically during program start
      \item[Tests] standalone callable performing runtime checks
    \end{description}

    \vspace{10pt}
    \begin{flushright}
      \includegraphics[width=.6\textwidth]{img/qr}
    \end{flushright}
  \end{column}
 \end{columns}

\end{frame}
\note{
  \begin{itemize}
    \item
    QR Code leads to the talk (clickable links)

    \item \textbf{Suites}: \textbf{automatically executed}, generate names, setup tests, nest
    \item \textbf{Name}: just a string, non-unique supported (but confusing)
    \item \textbf{Tests}: might be executed (runtime), data must not be dangling
    \item Special assignment for suites necessary at namespace scope.
    \item Both suites and tests are usually lambdas, but functions also work.
  \end{itemize}
}

\begin{frame}[fragile]
 \frametitle{Runtime Configurable}

 \begin{columns}[T,onlytextwidth]
  \begin{column}{.45\textwidth}
    \begin{console}[gobble = 6]
      \MG{>} ./demo
    \end{console}
    \includeConsole{Intro}
  \end{column}
  \begin{column}{.45\textwidth}
    \begin{console}[gobble = 6]
      \MG{>} ./demo \textminus{}\textminus{}list
    \end{console}
    \includeConsole{Intro-listing}
  \end{column}
 \end{columns}

 \begin{itemize}
  \item Test selection
  \item Controlling parallelism
  \item Further options with \code{\textminus{}\textminus{}help}
 \end{itemize}
\end{frame}
\note{
  \begin{itemize}
    \item
    Executes \textbf{all tests} by default

    \item
    Actual output shown (i.e.\ generated during build); in particular \textbf{all demos compile}.

    \item
    Can already see that tests are \textbf{executed concurrently}.

    \item
    Many execution details \textbf{runtime configurable}.
  \end{itemize}
}

\section{Expectations}
\begin{frame}[fragile]
  \frametitle{Expectations}

  \begin{columns}[onlytextwidth]
    \begin{column}{.5\textwidth}
      \includeDemo{ExpectIntro.cpp}
    \end{column}

    \begin{column}{.5\textwidth}
      \includeConsole{ExpectIntro}
    \end{column}
  \end{columns}
\end{frame}
\note{
  \begin{itemize}
    \item Sad truth: no introspection without a little help
    \item Second reason for using a testing framework
  \end{itemize}
}

\begin{frame}[fragile]
  \frametitle{Expectations with Expression Introspection}

  \begin{itemize}
    \item
    \code{ct::lift} wraps anything for introspection
    \item
    Many UDLs for lifting scalar and string types
  \end{itemize}

  \begin{columns}[onlytextwidth]
    \begin{column}{.5\textwidth}
      \includeDemo{ExpectIntrospection.cpp}
    \end{column}

    \begin{column}{.5\textwidth}
      \includeConsole{ExpectIntrospection}
    \end{column}
  \end{columns}

  \begin{itemize}
    \item
    Implemented with \emph{\href{https://en.wikipedia.org/wiki/Expression_templates}{Expression Templates}}
    [\href{https://web.archive.org/web/20050210090012/http://osl.iu.edu/~tveldhui/papers/Expression-Templates/exprtmpl.html}{T.\ Veldhuizen 1995},
    \href{https://www.amazon.de/-/en/David-Vandervoorde/dp/0201734842}{D.\ Vandevoorde 2002}]
  \end{itemize}
\end{frame}
\note{
  \begin{itemize}
    \item Expression tee as paramter (object)
    \item lifted value usable exactly as the original value.
  \end{itemize}
  \textbf{Meta}
  \begin{itemize}
    \item
    How does that work?

    \item
    Dive into \textbf{details now}, return to other \textbf{features later}.
  \end{itemize}
}

\subsection{Basic Attempt}
\begin{frame}[fragile]
  \frametitle{Expression Introspection Attempt (1/4)}
  \vspace{-4pt}
  \begin{columns}[T]
    \begin{column}{.5\textwidth}
      \includeDemo[%
        linerange = {6-29},%
        firstnumber = 6,%
        aboveskip = 0pt,%
        belowskip = 0pt,%
      ]{ExpectAttempt.cpp}
    \end{column}

    \begin{column}{.5\textwidth}
      \includeDemo[%
        linerange = {31-42},%
        firstnumber = 31,%
        aboveskip = 0pt,%
        belowskip = 0pt,%
      ]{ExpectAttempt.cpp}
    \end{column}
  \end{columns}
\end{frame}
\note{
  \begin{itemize}
    \item
    Conversion to bool for \code{static\_assert}
  \end{itemize}
}

\begin{frame}[fragile]
  \frametitle{Expression Introspection Attempt (2/4)}
  \vspace{-4pt}
  \begin{columns}[T]
    \begin{column}{.5\textwidth}
      \includeDemo[%
        linerange = {46-69},%
        firstnumber = 46,%
        aboveskip = 0pt,%
        belowskip = 0pt,%
      ]{ExpectAttempt.cpp}
    \end{column}

    \begin{column}{.5\textwidth}
      \includeDemo[%
        linerange = {71-80},%
        firstnumber = 71,%
        aboveskip = 0pt,%
        belowskip = 0pt,%
      ]{ExpectAttempt.cpp}
    \end{column}
  \end{columns}
\end{frame}
\note{
  \begin{itemize}
    \item \code{operator==}
    \begin{itemize}
      \item \code{lift} for convenience (with deduction guides)
      \item \code{requires} specs omitted
      (at least one \code{Expression}, \code{Lhs} and \code{Rhs} actually support \code{==})
    \end{itemize}
  \end{itemize}
}

\begin{frame}[fragile]
  \frametitle{Expression Introspection Attempt (3/4)}

  \begin{columns}[onlytextwidth]
    \begin{column}{.5\textwidth}
      \includeDemo[%
        linerange = {82-90},%
        firstnumber = 82,%
      ]{ExpectAttempt.cpp}
      \includeDemo[%
        linerange = {1-3,20},%
        firstnumber = 94,%
      ]{ExpectAttemptMain.h}
    \end{column}

    \begin{column}{.45\textwidth}
      \includeConsole[lastline=1]{ExpectAttempt}
      \begin{figure}
        \centering
        \includeDrawIo[width=.35\textwidth]{attempt0}
      \end{figure}
    \end{column}
  \end{columns}
\end{frame}
\note{
  \begin{itemize}
    \item Question to audience: Why could this be \textbf{problematic}?
    \item
      Using \code{t} \textbf{twice} in \code{expect} could be considered a \textbf{code smell}.
      It's actually even more than that.
    \item
      Can easily be extended to many other operators.
  \end{itemize}
}

\begin{frame}[fragile]
  \frametitle{Required Operators}
  \begin{block}{Overloaded Operators}
    \begin{tabular}{ l l }
      \textbf{Unary} &
      \code{not},
      \code{+},
      \code{-},
      \code{*} (dereference)
      \\

      \textbf{Arithmetic}
      &
      \code{+},
      \code{-},
      \code{*},
      \code{/},
      \code{\%}
      \\

      \textbf{Bitwise}
      &
      \code{\&},
      \code{|},
      \code{\textasciitilde{}},
      \code{\textasciicircum{}}
      \\

      \textbf{Comparison}
      &
      \code{<},
      \code{<=},
      \code{>},
      \code{>=},
      \code{==},
      \code{!=}
      \\

      \textbf{Special}
      &
      \code{,} (comma),
      \code{and},
      \code{or}
      \\
    \end{tabular}
  \end{block}
\end{frame}
\note{
  \begin{itemize}
    \item
    22 operators in total with \textbf{\code{const} parameters}

    \item
    Special \textbf{text syntax} for \code{not}, \code{and}, and \code{or}.

    \item
    Implementations for \code{and} and \code{or} more \textbf{involved}.

    \item
    \textbf{Hard} to get right; start with \textbf{simple attempt}.
  \end{itemize}
}

\begin{frame}[fragile]
  \frametitle{Expression Introspection Attempt (4/4)}
  \vspace{-4pt}
  \begin{columns}[T]
    \begin{column}{.5\textwidth}
      \includeDemo[%
        linerange = {92-115},%
        firstnumber = 92,%
        aboveskip = 0pt,%
        belowskip = 0pt,%
      ]{ExpectAttempt.cpp}
    \end{column}

    \begin{column}{.5\textwidth}
      \includeDemo[%
        linerange = {1-7,20},%
        firstnumber = 184,%
        aboveskip = 0pt,%
      ]{ExpectAttemptMain.h}
      \visible<2 | handout:1> {
        \includeConsole[lastline=2]{ExpectAttempt}
      }
      \vspace{-5pt}
      \begin{center}
        \includeDrawIo[width=.38\textwidth]{attempt1}
      \end{center}
    \end{column}
  \end{columns}
\end{frame}
\note{
  \begin{itemize}
    \item
    Custom \code{operator or} never short circuits,
    but rather works as function call.
  \end{itemize}
  % https://stackoverflow.com/questions/35918735/short-circuit-evaluation-on-c-logical-and-operators#comment59495892_35918735

  \textbf{Problem Analysis}
  \begin{itemize}
    \item
    Output needs to be short circuit aware

    \item
    Re-evaluating \code{m\_lhs.value} infeasible

    \item
    Need caching

    \item
    Discover my next attempt

    \item
    ($\Rightarrow$ Naturally can't deal with rvalue refs transparently.)
  \end{itemize}
}

\subsection{Caching Attempt}
\begin{frame}
  \frametitle{Expression Introspection Cached Attempt (1/2)}

  \begin{columns}[T]
    \begin{column}{.5\textwidth}
      \includeDemo[%
        linerange = {92-112},%
        firstnumber = 92,%
      ]{ExpectCached.cpp}
    \end{column}
    \begin{column}{.5\textwidth}
      \includeDemo[%
        linerange = {113-132},%
        firstnumber = 113,%
      ]{ExpectCached.cpp}
    \end{column}
  \end{columns}
\end{frame}
\note{}

\begin{frame}
  \frametitle{Expression Introspection Cached Attempt (2/2)}

  \begin{columns}[onlytextwidth]
    \begin{column}{.5\textwidth}
      \begin{overlayarea}{\textwidth}{.6\textheight}
        \only<1 | handout:0>{
          \includeDemo[%
            linerange = {1-7,20},%
            firstnumber = 184,%
          ]{ExpectAttemptMain.h}
        }
        \only<2->{
          \includeDemo[%
            linerange = {1-12,20},%
            firstnumber = 184,%
          ]{ExpectAttemptMain.h}
        }
      \end{overlayarea}

      \begin{overlayarea}{\textwidth}{.2\textheight}
        \only<1 | handout:0>{
          \includeConsole[lastline=2]{ExpectCached}
        }
        \only<2->{
          \includeConsole{ExpectCached}
        }
      \end{overlayarea}
    \end{column}

    \begin{column}{.5\textwidth}
      \only<2->{
          \begin{figure}
            \centering
            \includeDrawIo[width=.9\textwidth]{cached}
          \end{figure}
      }
      \begin{overlayarea}{\textwidth}{.2\textheight}
        \only<3->{
          \begin{itemize}
            \item
            May only cache upon \code{value()} call

            \item
            Requires \code{mutable} members;\\
            not \code{constexpr}.
          \end{itemize}
        }
      \end{overlayarea}
    \end{column}
  \end{columns}
\end{frame}
\note{
  \begin{itemize}
    \item
    Seemed like a dead end
  \end{itemize}
}

\begin{frame}[standout]
    \begin{quote}
      All problems in computer science can be solved \\
      by another level of indirection.
    \end{quote}
    \begin{flushright}
      David J.\ Wheeler
    \end{flushright}
\end{frame}
\note{
  \begin{itemize}
    \item
    "Fundamental theorem of software engineering"

    \item
    \textbf{separate} expression and evaluation

    \item
    Evaluation contains cached (sub-)values
  \end{itemize}
}

\subsection{Final Version}
\begin{frame}[fragile]
  \frametitle{Expression Evaluation with Introspection (1/3)}

  \parbox[c][.92\textheight][c]{\textwidth}{
    \includeDemo[%
      linerange = {176-197},%
      firstnumber = 176,%
      aboveskip = 0pt,%
      belowskip = 0pt,%
    ]{ExpectEvaluation.cpp}
  }
\end{frame}
\note{
  \begin{itemize}
    \item
    Split API
  \end{itemize}
}

\begin{frame}
  \frametitle{Expression Evaluation with Introspection (2/3)}

  \begin{columns}[T,onlytextwidth]
    \begin{column}{.5\textwidth}
      \includeDemo[%
        linerange = {137-159},%
        firstnumber = 137,%
      ]{ExpectEvaluation.cpp}
    \end{column}
    \begin{column}{.5\textwidth}
      \includeDemo[%
        linerange = {160-174},%
        firstnumber = 160,%
      ]{ExpectEvaluation.cpp}
    \end{column}
  \end{columns}
\end{frame}
\note{
  \begin{itemize}
    \item
    Carry over idea to all other operators

    \item
    Reduce boilerplate with a CRTP base
  \end{itemize}
}

\begin{frame}
  \frametitle{Expression Evaluation with Introspection (3/3)}

  \begin{columns}[T,onlytextwidth]
    \begin{column}{.5\textwidth}
      \includeDemo[%
        linerange = {61-71},%
        firstnumber = 61,%
      ]{ExpectEvaluation.cpp}
      \includeDemo[%
        linerange = {1-12,20},%
        firstnumber = 184,%
      ]{ExpectAttemptMain.h}
    \end{column}
    \begin{column}{.5\textwidth}
      \includeConsole{ExpectEvaluation}
      \begin{center}
        \includeDrawIo[%
          width=\textwidth,%
        ]{evaluation}
      \end{center}
    \end{column}
  \end{columns}
\end{frame}
\note{
  \begin{itemize}
    \item
    \code{ct::expect} simplified: only accepts \code{Expression}s here.

    \item
    3 cases of evaluating \code{or}s: evaluating 0, 1, 2 children.
  \end{itemize}
}

\subsection{Limitations}
\begin{frame}
  \frametitle{Limitations of \texttt{ct::lift}}

  \begin{itemize}
    \item
    Ensure that operators have at least one \code{ct::Expression} operand, e.g.\ \code{ct::lift}.

    \item
    Function calls are not lazy by default
    \includeDemo[%
      linerange = {15-24},%
      firstnumber = 15,%
    ]{Limitations.cpp}
  \end{itemize}
\end{frame}
\note{}

\section{Advanced Expectations}
\begin{frame}[fragile]
  \frametitle{Advanced Expectations}

  \begin{columns}[T,onlytextwidth]
    \begin{column}{.56\textwidth}
      \setlength{\leftmargini}{0.0cm}
      \begin{itemize}
        \item
          Debug output
          \includeDemo[%
            linerange = {22-22},%
            firstnumber = 22,%
            aboveskip = 4pt,%
            belowskip = 7pt,%
          ]{AdvancedExpectations.cpp}
        \item
          Assertions: \code{asserted} / \code{asserted\_if}
          \includeDemo[%
            linerange = {26-27},%
            firstnumber = 26,%
            aboveskip = 4pt,%
            belowskip = 7pt,%
          ]{AdvancedExpectations.cpp}
        \item
          Flakyness: \code{flaky} / \code{flaky\_if}
          \includeDemo[%
            linerange = {33-33},%
            firstnumber = 33,%
            aboveskip = 4pt,%
            belowskip = 7pt,%
          ]{AdvancedExpectations.cpp}
        \item
          Exceptions: \code{throws} / \code{throws<Exception>}
          \includeDemo[%
            linerange = {39-39},%
            firstnumber = 39,%
            aboveskip = 4pt,%
            belowskip = 7pt,%
          ]{AdvancedExpectations.cpp}
        \item
          Fatal Assertions: \code{aborts} / \code{debug\_aborts}
          \includeDemo[%
            linerange = {49-51},%
            firstnumber = 49,%
            aboveskip = 4pt,%
            belowskip = 7pt,%2
          ]{AdvancedExpectations.cpp}
      \end{itemize}
    \end{column}
    \begin{column}{.46\textwidth}
      \includeConsole{AdvancedExpectations}
    \end{column}
  \end{columns}
\end{frame}
\note{
  \begin{itemize}
    \item
    \code{noexcept} via \code{not throws}; all part of the expression realm.

    \item
    Abort / throw checks \textbf{lazy}.
  \end{itemize}
}

\section{Distance Checks}
\begin{frame}[fragile]
  \frametitle{Advanced Distance Checks}

  \vspace{-4pt}
  \begin{columns}[T,onlytextwidth]
    \begin{column}{.55\textwidth}
      \includeDemo[%
        linerange = {17-27},%
        firstnumber = 17,%
      ]{Distance.cpp}
    \end{column}
    \begin{column}{.45\textwidth}
      \begin{itemize}
        \item
          CPO \code{ct::utils::norm}

        \item
          Absolute distance
          \[
            \Delta_{\text{abs}}\left(x, y\right)
            =
            \left\lVert x - y \right\rVert
          \]

        \item Relative distance
          \[
            \Delta_{\text{rel}}\left(x, y\right)
            =
            \frac{
              \left\lVert x - y \right\rVert
            }{
              \max\left(
                \left\lVert x \right\rVert,
                \left\lVert y \right\rVert
              \right)
            }
          \]
      \end{itemize}
    \end{column}
  \end{columns}
  \vspace{4pt}

  \includeConsole{Distance}
\end{frame}
\note{
  \begin{itemize}
    \item
    \code{ct::tolerance(absolute, relative)}, defaults: \code{Limits::epsilon()}

    \item
    \code{ct::utils::norm} different from \code{std::norm}: returns \code{T} for most types, extensible via ADL and member functions.

    \item
    (ULP (units in the last place; float interpreted as int) not available (yet, since hard without UB).)
  \end{itemize}
}

\begin{frame}
  \frametitle{Conversion Tools}

  \textbf{\href{%
    https://github.com/clean-test/migration%
  }{
    \includegraphics[width=.65\baselineskip]{img/github}
    clean-test/migration
  }
  }

  \begin{itemize}
    \item
    Hackable converter for existing tests

    \item
    \code{ct::lift} expectations for introspection
  \end{itemize}

  \vspace{6pt}
  \begin{columns}[T,onlytextwidth]
    \begin{column}{.5\textwidth}
      \emph{Original}
      \includeDemo{ConversionBefore.cpp}
    \end{column}
    \begin{column}{.5\textwidth}
      \emph{Converted}
      \includeDemo{ConversionAfter.cpp}
    \end{column}
  \end{columns}
\end{frame}
\note{
  \begin{itemize}
    \item
    Written in python

    \item
    Attempts to lift the tree appropriately (heuristic, works well)

    \item
    Currently boost only

    \item
    Shown example \textbf{without modifications}

    \item
    Supports \textbf{all test features} seen here, e.g.\ distance checks

    \item
    Output not perfect in general: might need \textbf{minor manual adaptations} to compile
  \end{itemize}
}

\section{Concurrency}
\begin{frame}[fragile]
  \frametitle{Concurrent Tests}

  \begin{columns}[onlytextwidth]
    \setlength{\leftmargini}{-5pt}
    \begin{column}{.5\textwidth}
      \begin{itemize}
        \item
        Test-case attribution managed via \code{ct::Observer}

        \item
        Automatic for single threaded tests

        \item
        Propagate \code{ct::Observer} for advanced parallel tests

        \item
        \code{ct::expect} thread safe
      \end{itemize}
    \end{column}
    \begin{column}{.5\textwidth}
      \includeDemo[
        linerange = {20},
        firstnumber = 20,
      ]{Observer.cpp}
    \end{column}
  \end{columns}

  \includeConsole{Observer}
\end{frame}
\note{
  \begin{itemize}
    \item
    Two more \textbf{high level features}:

    \item
    Step back: single threaded execution makes \textbf{linking / connecting} failures to tests trivial.

    \item
    2 options for propagation:
    1st \code{ct::expect} parameter or
    via \code{ct::ObservationSetup}.

    \item
    \textbf{Warning} identifies all affected and \RD{executed} locations.
    Note that this also applies to successful expectations.
  \end{itemize}
}

\section{Data Driven}
\begin{frame}[fragile]
  \frametitle{Data Driven Tests}

  \begin{columns}[onlytextwidth]
    \begin{column}{.55\textwidth}
      \includeDemo[%
        linerange = {15-31},%
        firstnumber = 15,%
      ]{DataDriven.cpp}
    \end{column}
    \begin{column}{.47\textwidth}
      \includeConsole{DataDriven}
    \end{column}
  \end{columns}
\end{frame}
\note{
  \begin{itemize}
    \item
    Personal favorite for verifying invariants
    \item
    Two ways of passing data: parameter and pipe.

    \item
    Can be combined with \code{Observer}s.

    \item
    Copes with rvalue reference data.
  \end{itemize}
}

\begin{frame}
  \frametitle{Wrap-up}

  \vspace{8pt}
  \begin{exampleblock}{Summary: Clean Test}
    \begin{itemize}
      \item is a modern, versatile and yet simple to use testing framework.
      \item supports short-circuit expression introspection without macros.
      \item provides various productivity features.
      \item is built for parallel tests and test execution.
    \end{itemize}
  \end{exampleblock}

  \begin{columns}[onlytextwidth]
    \begin{column}{.48\textwidth}
      \begin{block}{Future Work}
        \parbox[c][.35\textheight][c]{\textwidth}{
          \begin{itemize}
            \item Optimize test scheduling
            \item Lazy data providers
            \item Further migration utilities
            \item Convenience (e.g.\ for ranges)
          \end{itemize}
        }
      \end{block}
    \end{column}

    \begin{column}{.48\textwidth}
      \begin{block}{References}
        \parbox[c][.35\textheight][c]{\textwidth}{
          \begin{tabular}{ @{}l l@{} }
            Clean Test &
            \href{%
              https://github.com/clean-test/clean-test
            }{%
              \includegraphics[height=10pt]{img/github} clean-test/clean-test
            } \\

            Migrate&
            \href{%
              https://github.com/clean-test/migration
            }{%
              \includegraphics[height=10pt]{img/github} clean-test/migration
            } \\

            This Talk&
            \href{%
              https://github.com/clean-test/talk/tree/main/2022-meeting-cpp
            }{%
              \includegraphics[height=10pt]{img/github} clean-test/talk
            } \\

            Philipp &
            \href{https://github.com/m8mble/}{\includegraphics[height=10pt]{img/github}}
            \href{https://stackoverflow.com/users/2747160/m8mble}{\includegraphics[height=10pt]{img/stackoverflow}}
            \href{https://twitter.com/m8mble}{\includegraphics[height=10pt]{img/twitter}}
            \href{https://mastodon.social/@m8mble}{\includegraphics[height=10pt]{img/mastodon}}
            \href{https://www.reddit.com/user/m8mble}{\includegraphics[height=10pt]{img/reddit}}
            \texttt{@m8mble}
          \end{tabular}\\

          Comments, issues and PRs welcome.
        }
      \end{block}
    \end{column}
  \end{columns}
\end{frame}
\note{
  \begin{itemize}
    \item
    Make a 1.0 release

    \item
    Appreciate any sort of feedback

    \item
    Scheduling: Ask for a certain amount of (usable) CPUs.
  \end{itemize}
}

%%%%%%%%%%%%%%%%%%%%%%%%%%%%%%%%%%%%%%%%%%%%%%%%%%%%%%%%%%%%%%%%%%%%%%%%%%%%%%%%%%%%%%%%%%%%%%%%%

\appendix

\begin{frame}
  \frametitle{Asio Thread Pool}

  \vspace{-4pt}
  \begin{columns}[T]
    \begin{column}{.5\textwidth}
      \includeDemo[%
        linerange = {15-37},%
        firstnumber = 11,%
        aboveskip = 0pt,%
        belowskip = 0pt,%
      ]{Asio.cpp}
    \end{column}

    \begin{column}{.5\textwidth}
      \includeDemo[%
        linerange = {40},%
        firstnumber = 35,%
        aboveskip = 0pt,%
        belowskip = 0pt,%
      ]{Asio.cpp}

      \vspace{10pt}
      \includeConsole{Asio}
    \end{column}
  \end{columns}
\end{frame}

\end{document}

